\documentclass[11pt,a4paper,DIV=12,numbers=noenddot,bib=totoc]{scrreprt}

% Helvetica font
%\usepackage[scaled]{helvet}
%\renewcommand{\familydefault}{\sfdefault}

% Set line spacing to match previous 11pt layout
\usepackage{setspace}
\setstretch{1.5}  % increases line spacing to preserve page coverage

% Page margins
\usepackage[left=2.5cm,right=2.5cm,top=1.5cm,bottom=1.5cm,includeheadfoot]{geometry}

% Optional: keep zero paragraph spacing if desired
\setlength{\parindent}{15pt}
\setlength{\parskip}{0pt}



%footnote in caption
\usepackage{ftnxtra}
\usepackage{fnpos}

%Einbinden von pdf
\usepackage{pdfpages}

%Einbinden von Python code
\usepackage{minted}

%Tabellen
\usepackage[table,xcdraw]{xcolor}

\usepackage{indentfirst}

%Package für Tabelle bei MatMeth Hardened properties
\usepackage{graphicx}
\usepackage{booktabs}
\usepackage{tabularx}

%Tikz - Package für Graphen
\usepackage{tikz, pgfplots, pgf}

%Mathematik
\usepackage{amsmath,amssymb}
\usepackage{mathtools}

%Subfigures
\usepackage{subfigure}
\usepackage{subcaption} % this is essential


%Textumflossene Bilder
\usepackage{wrapfig}

%Aufzählungen individuell nummereieren
\usepackage{enumitem}

%Verbesserung für Floatobjekte
\usepackage{float}

%Tex zur Ausgabe von Floatobjekten zwingen mittels \FloatBarrier
\usepackage{placeins}

% Float objekten andere caption zuweisen
\usepackage{capt-of}

%Paket für Promillezeichen
\usepackage{textcomp}

% Farbbox
\usepackage[most]{tcolorbox}
% Define ETH Blue
\definecolor{ETHgreen}{HTML}{627313}
\usepackage[table]{xcolor} % load xcolor with table support
\definecolor{ETHblue}{HTML}{003399}
\definecolor{ETHgreen}{HTML}{627313}

%Tikz 
\usetikzlibrary{arrows.meta, positioning, shapes.geometric, calc}
\usetikzlibrary{decorations.pathreplacing}
%\usepackage{amssymb}


%Literaturverzeichnis
%\usepackage[sorting=none, style=numeric]{biblatex}
%\addbibresource{Bib.bib}
\usepackage[numbers, square, sort&compress]{natbib}
\bibliographystyle{unsrtnat}
\setlength{\bibsep}{0.2ex}  % default ~1.5ex, adjust as needed


%Verzeichnisse zulassen
%\usepackage[nonumberlist, acronym, section]{glossaries}
%Farben
\usepackage{color}
\definecolor{lg}{gray}{0.95}
\definecolor{dg}{gray}{0.35}
\definecolor{gr}{HTML}{007F0E}
\definecolor{bl}{HTML}{00137F}
\definecolor{em}{HTML}{088A68}

%verweise im Dokument
\usepackage{xurl}
\usepackage{hyperref}
\hypersetup{colorlinks=true,
    linkcolor = black,
    citecolor = black,
    urlcolor  = black} 
\urlstyle{same}

\usepackage{multirow}

%probetext
\usepackage{blindtext}

%Matlab code einbinden
\usepackage{matlab-prettifier}



% Seitenstil
\usepackage[]{scrlayer-scrpage}
\pagestyle{scrheadings}

\clearpairofpagestyles
%Kopfzeile
%\ihead[\textsf{Improving the analysis of potential mapping}]
%\setheadsepline{0.5pt}

%Fußzeile
%\setfootsepline{0.5pt}
\ifoot[]{}
\cfoot[]{}
\ofoot[\pagemark]{\pagemark}

%schriftart
%\renewcommand{\rmdefault}{\sfdefault}

%Gesamtseitenzahl
\usepackage{lastpage}

%Verzeichnisse anpassen
\usepackage{tocbasic}
%\usetocstyle{classic}depricaded
%\settocfeature[toc][0]{entryvskip}{0.9em plus .2pt}

\makeatletter
%paragraph umdefinieren
%\renewcommand\paragraph{\@startsection{paragraph}{4}{\z@}%
%  {-3.25ex\@plus -1ex \@minus -.2ex}%
%  {1.5ex \@plus .2ex}%
%	{\normalfont\normalsize\bfseries}}
	
%römische Ziffer
\newcommand{\rom}[1]{\textsc{\@roman{#1.}}}

\makeatother

% Global paragraph formatting
\setlength{\parindent}{0pt}  % Indent new paragraphs
\setlength{\parskip}{0pt}     % No extra vertical space between paragraphs


%floating einstellungen für die anteile an float und text auf einer seite
\renewcommand{\topfraction}{1.0}
\renewcommand{\bottomfraction}{1.0}
\renewcommand{\textfraction}{0.0}

% Section spacing adjustments
\RedeclareSectionCommand[
  beforeskip=2ex plus 1ex minus 0.5ex,  % space before section
  afterskip=0.75ex                            % space after section
]{section}

\RedeclareSectionCommand[
  beforeskip=1.5ex plus 0.5ex minus 0.5ex, % space before subsection
  afterskip=0.5ex                           % space after subsection
]{subsection}

\RedeclareSectionCommand[
  beforeskip=1ex plus 0.3ex minus 0.3ex,   % space before subsubsection
  afterskip=0.5ex                            % space after subsubsection
]{subsubsection}





\begin{document}
\section{Kapitelhierarchie und Dokumentorgnisation}\label{sec:Struktur}
Es gilt: Chapter>Section>Subsection>Subsubsection>Paragraph, wobei bei einem kurzen Bericht selten auf die 3. Ebene, fast nie auf die 4. Ebene runter gegangen wird
%%%
\subsection{Subsection-Beispiel}\label{sec:Struktur-Subsection}
Das ist eine Subsection.
%%%%
\subsubsection{Subsubsection-Beispiel}\label{sec:Struktur-Subsection-Subsubsection}
Das Ist eine Subsubsection. Es ist die 4. Hierarchieebene, weshalb auf Darstellung der Nummer verzichtet wird. In einem kurzen Bericht sollte man nicht so tief gehen.
%%%%
\paragraph{Paragraph-Beispiel} kann referenziert werden und notfalls auch im Inhaltsverzeichnis aufgenommen werden (tocdepth=5). Es gibt keinen Zähler für Paragraphen.
%%%%%%%%%%%%
\section{Schriftgrössen definieren} \label{sec:Schriften}
Schriftgrössen sind. {\tiny Winzig} {\footnotesize Fussnote} {\small Kleiner} {\normalsize Normal} {\Large Grösser} {\huge riesig} {\Huge Gigantisch}. Soll Text betont werden, wird er über \textit{kursiv gesetzt}. Fetter Text ist \textbf{mit textbf.} Kommentiert wird mit \%. 

Schriften sollten im Header oder Stylefile geändert werden, da wir Inhalte und Formatierung in Latex trennen möchten. Grundsätzlich gibt es keine Vorgaben, nur sollten im Dokument verwendete Schriften durchgängig sein.
%%%%%%%%
\section{Formeln-Beispiele und Umgebungen}\label{sec:Formeln}
Formeln können im Text mit Dollarzeichen über $5=2>a\cdot \alpha$ ohne Zähler eingebettet werden. Formeln haben einen besonderen Zeichensatz, weshalb man Formelsymbole immer zwischen \$ \$ setzen sollte. Eine Formel kann natürlich auch abgesetzt werden wie die hier:
\begin{equation}\label{eq:Beispielformel}
P_\text{3.O} = a_0 +a_1\omega + a_2 \omega^2 +a_3\omega^3.
\end{equation}
Diese werden durchnummeriert (ausser equation*) und können über das Label referenziert werden und zwar als Gl.~\ref{eq:Beispielformel}. Da hyperref verwendet wird erzeugt dies im Pdf-Dokument einen Link. Achtung immer daran denken dass Formeln Teil des Satzes sind, also mit Punkt und Komma. Formeln sehen in Tex einfach super aus. Die Formelsyntax ist in Latex allerdings gewöhnungsbedürftig. Man muss entweder nachschauen (z.B. Ref.~\cite{wiki_book,BeuermannSascha2002EvlP}) oder kann auch andere Formeleditoren verwenden und dann in Latex konvertieren (z.B. Mathtype). Mehrzeilige Formeln benötigen die eqnarray Umgebung.
%%%%%%%%%%%%
\section{Tabellen-Beispiele und Umgebungen}\label{sec:Tabellen}
Tabellen sind in Latex etwas gewöhnungsbedürftig. Es gibt unterschiedliche Tabellenumgebungen (tabular, longtable, table, tabbing, etc.). Tabellen sind wie Abbildungen Floating-objekte wie auch Bilder, werden also vom Algorithmus im Dokument platziert. Man macht Spaltendefinitionen, die Beschriftung (caption) und dann den Inhalt. Spaltendefinitionen sind l=links; c=zentriert; r=rechts, p\{Spaltenbreite\}, $\|=$ senkrechte Tabellenlinie.  Ein Beispiel:%Verbatim  erlaubt es code zu drucken ohne dass es interpretiert wird.
\begin{verbatim} 
\begin{table}[H]
\centering
   \caption{Beispiel einer einfachen Tabelle.}   \label{tab:Beispiel}
\begin{tabular}{|p{60mm}|l|}
\hline
links oben               &     rechts oben  \\ \hline \hline
links unten   & rechts unten   \\ \hline
\end{tabular}
\end{table}
\end{verbatim}
\begin{table}[H]
\centering
   \caption{Beispiel einer einfachen Tabelle.}   \label{tab:Beispiel}
\begin{tabular}{|p{60mm}|l|}
\hline
links oben                 &     rechts oben  \\ \hline \hline
links unten   & rechts unten   \\ \hline
\end{tabular}
\end{table}
Es gibt auch kleine Hilfsprogramme (Tabelleneditoren), die die Erstellung richtig grosser Tabellen erheblich erleichtern können als kostenlose Programme, zur Konversion ganzer Excell-sheets oder auch als online converter (z.B. \url{https://www.tablesgenerator.com}), aus denen man dann nur noch die \LaTeX -Passage ins Dokument kopiert. Die Tabellenbeschriftung muss {\scshape\Large immer über} die Tabelle (s. Tab.~\ref{tab:Beispiel}) und bestehen wie bei Abbildungen aus einem vollständigen Satz.
%%%%%%%%%
\section{Abbildungen-Beispiele und Umgebungen}\label{sec:Abbildungen}
Abbildungen sind wie Tabellen Floating-objekte, die nach Formatvorgaben verteilt werden. Anweisungen sind h=hier; t=top Seite; b=Bottom Seite; p=Platzierung auf eigener Seite. Mit [H] kann Platzierung erzwungen werden, was aber nicht von Vorteil ist. Die Grössenangabe wird über width oder height gemacht. Es gibt auch scale, angle. Angaben sind in pt, mm, cm, oder Anteile der Zeilenbreite (textwidth) oder Seitenbreite. Bilder so klein wie möglich halten und wenn möglich PDF Vektorgrafiken verwenden. PDF, JPG und PNG Formate sind problemlos. Ein erstes Beispiel ist in Abb.~\ref{fig:Beispielbild1} gezeigt.
 \begin{figure}[htb]
	\begin{center}
	\includegraphics[width=0.85\textwidth]{IMAGES/writing_tutorial_fig3.jpg}
    \caption{\label{fig:Beispielbild1}Beispielbild unterschiedlicher Betonmischungen C1 und C2. (a) C1 nach der Mischung, (b) C2 ohne Verflüssiger und (c) C2 mit Verflüssiger.}
  	\end{center}
\end{figure}
%%%%%%%%%
\section{Referenzen-Beispiele}\label{sec:referenzen}
Alle referenzierten Objekte (Formeln, Tabellen, Kapitel, Abbildungen) haben eigene Zähler. Überall sollte ein label gesetzt sein. Sinnvoll sind label, die was aussagen z.B. ein Label für Gleichungen als label\{eq:Mitternachtsformel\}, für Abbildungen label\{fig:Beispielbild1\}, Tabellen \{tab:Materialdaten\} u.s.w.. Wenn man dann referenziert dann erfolgt das mit dem Befehlt ref{eq:Mitternachtformel}. Ich würde immer gleich jedes definierte Objekt mit einem Label versehen, das macht das Arbeiten leichter. Bisher haben wir Gleichung~\ref{eq:Beispielformel}, Tabelle~\ref{tab:Beispiel}, Abbildungen~\ref{fig:Beispielbild1}-\ref{fig:Beispielbild3} in Kapitel~\ref{chap:mm} definiert. Es ist das hyperref Paket im header gesetzt, d.h. im fertigen PDF Dokument sind alle Verweise, Kapitel in der Inhaltsangabe und Quellenangaben als Hyperlinks eingebettet (farbig markiert).
%%%%%%%%%%%
\section{Aufzählungen-Beispiele und Umgebungen}\label{sec:itemize}
Aufzählungen sind ideal um komplizierte, umfangreiche Sachverhalte zu strukturieren. Sie können als bullit items oder Aufzählungen sein: Ein Beispiel für Aufzälungen:
\begin{itemize}
\item Erstes Item
\item[$\Omega$] zweites Item. spezielles Item zwischen eckigen Klammern.
\end{itemize}
Aufzählungen können auch geschachtelt sein in allen möglichen Kombinationen:
\begin{enumerate}
  \item Erster Punkt
  \begin{enumerate}
     \item Erster Unterpunkt des ersten Punkten
     \item Zweiter Unterpunkt des ersten Punkten
  \end{enumerate}
  \item Zweiter Punkt
  \begin{itemize}
     \item Erstes Item
     \item[$\Omega$] zweites Item. spezielles Item zwischen Klammern. \[\].
  \end{itemize}
\end{enumerate}

Hier \glqq Weisheiten\grqq~ zum Bericht schreiben als Aufzählung:
\begin{itemize}
    \item Sei dir bewusst, dass die Erstellung eines Berichtes mehrere Stunden braucht, das ist bei der Anzahl Kreditpunkte der Lehrveranstaltung eingerechnet.
    \item Denk dran, Berichte sollen informativ sein und die Sache/das Problem steht im Vordergrund (kein Ich/Wir).
    \item Sei klar bei den Zielen des Berichtes.
    \item Verwende einfache, kurze Sätze, aber keine Emotionen oder Umgangssprache.
    \item Sei exakt und strukturiert beim Layout.
    \item KISS = keep it short \& simple; konzentriere dich auf das wesentliche und verwende auch nur wesentliches.
    \item Niemals Ergebnisse/Beobachtungen mit Interpretationen vermischen. Das sind unterschiedliche Kapitel.
    \item Sei knapp bei der Beschreibung bekannten und zitiere hier (formal korrekt). Investiere die Zeit lieber bei den Schlussfolgerungen.
    \item Sei deiner Arbeit gegenüber kritisch und versuche von anderen zu lernen, aber kopiere {\scshape\Large nie} von \grqq Freunden\grqq.
\end{itemize}
%%%%%%%%%%%%
\section{Quellenangaben}
Zitieren in \LaTeX ist sehr einfach. Es werden alles Quellen in einer *.bib Datei in Form von Bib-items gesammelt. Aus vielen Webseiten und anderen Programmen (wie Mendeley) lassen sich direkt bib-Einträge exportieren und kopieren. Ein Beispiel für eine Buch, eine Webseite und einen Artikel (mehr Beispiele in der {\em bibliography.bib} Datei):
\begin{verbatim}
@book{Weibull_McCool_2012,
title = "{Using the Weibull Distribution}",
author = "J.I.McCool",
year = "2012",
publisher = "Wiley"
}

@misc{wiki_book,
   author = "{Wikipedia contributors}",
   title = "Latex - Featured WikiBook",
   year = "2019",
   url = "https://en.wikibooks.org/wiki/LaTeX",
   note = "[Online; accessed 23-Octovber-2019]"
 }

@article{BeuermannSascha2002EvlP,
publisher = {Institut für Baumechanik und Numerische Mechanik, Uni Hannover; },
year = {2002},
title = {Erstellung von leistungsfähigen PDF-Dokumenten mit LATEX und den Paketen 
hyperref sowie thumbpdf},
edition = {Version 2.10},
author = {Beuermann, Sascha},
url        = {https://www.research-collection.ethz.ch/handle/20.500.11850/146361}
}
\end{verbatim}
Im Text zitiert wird über den Befehl {\em cite} die Quellen \cite{Weibull_McCool_2012,wiki_book,BeuermannSascha2002EvlP,Leigh-Phoenix-Beyerlein_2000}. Das Literaturverzeichnis wird automatisch erstellt. mit {\em citet} erhält man den Autor. Bsp: In \citet{BeuermannSascha2002EvlP} findet sich ....


Was bedeutet das alles? Hier erfolgt die Interpretation / Einordnung der Ergebnisse. Es wird erklärt was die Daten bedeuten und wie die Ergebnisse zu interpretieren sind. Man greift die Fragestellung der Einleitung wieder auf (ohne Wiederholung) und ordnet die Ergebnisse in das bisherige Wissen ein. Zwecks Lesefreundlichkeit können hierfür die relevanten Ergebnisse kurz aufgegriffen werden, eine ausführliche Wiederholung ist zu vermeiden. 

In der Diskussion werden die verschiedenen Ergebnisse auch verknüpft interpretiert, daher wird die Diskussion weniger in Unterkapitel unterteilt. Absätze können und sollen enthalten sein.

Als letzte(s) Unterkapitel der Diskussion finden sich Schlussfolgerungen/Fazit oder ein Ausblick/Empfehlungen, usw. wieder. Je nach Format können dies auch eigene Kapitel sein. 

 \begin{figure}[H]
	\begin{center}
	\includegraphics[width=0.8\textwidth]{IMAGES/writing_tutorial_fig2.png}
    \caption{\label{fig:Beispielbild3} Flussdiagramm des iterativen Arbeitsablaufs beim Abfassen von Berichten. (Quelle: \cite{1p_tutorial})}
  	\end{center}
\end{figure}
Schreiben ist ein iterativer Prozess, der mit der Konzeption anfängt, mit dem Aufbau, also dem anfertigen des Skeletts der Arbeit weitergeht und in Schritt 3 nach Abb.~ \ref{fig:Beispielbild3} mit der Ausformulierung weiter geht. An Ende steht das Korrekturlesen. Zudem sollten sich die Autoren fragen, ob sie die Checkliste in Abb.~\ref{fig:Beispielbild4} wirklich erfüllen.
 \begin{figure}[H]
	\begin{center}
	\includegraphics[width=0.9\textwidth]{IMAGES/writing_tutorial_fig4.pdf}
    \caption{\label{fig:Beispielbild4} Checkliste und Grundlage des Bewertungs
    schemas für Berichte. (Quelle: \cite{1p_tutorial})}
  	\end{center}
\end{figure}

\end{document}